\chapter{Tracciamento completo dei requisiti}
\label{app:requisiti}

In questa appendice sono raccolte le tabelle complete dei requisiti
(funzionali, qualitativi e di vincolo) con il relativo tracciamento sui casi
d'uso. La struttura del codice identificativo dei requisiti è descritta nella
sezione~\ref{sec:codice-requisiti}.

\begin{table}[h]
\caption{Tabella del tracciamento dei requisiti funzionali}
\label{tab:requisiti-funzionali}
\begin{tabularx}{\textwidth}{lXl}
\hline\hline
\textbf{Requisito} & \textbf{Descrizione} & \textbf{Use Case}\\
\hline
RFN-1 & Tracciamento del percorso completo di ogni richiesta HTTP, incluse le operazioni sul database & UC2, UC4 \\
\hline
RFN-2 & Identificazione delle operazioni CRUD tramite attributi di metodo e repository & UC2, UC4 \\
\hline
RFN-3 & Esposizione di metriche operative per i moduli di autenticazione, accesso ai dati e griglia & UC2 \\
\hline
RFN-4 & Correlazione tra log e trace tramite \texttt{TraceId} e \texttt{SpanId} & UC4 \\
\hline
RFN-5 & Attivazione/disattivazione della telemetria tramite configurazione \emph{appsettings} & UC3 \\
\hline
RFN-6 & Configurazione indipendente del campionamento per ciascuna sorgente di trace & UC3 \\
\hline
RFN-7 & Visualizzazione di trace e metriche raggruppate in dashboard Grafana & UC2 \\
\hline
RFN-8 & Raccolta dei dati da più servizi e installazioni, con filtraggio per cliente e servizio & UC2 \\
\hline
RFN-9 & Sviluppo di \emph{Claude skill} per la propagazione della telemetria ai moduli del framework & UC3 \\
\hline
RFN-10 & Arricchimento automatico degli span con identità utente (\texttt{enduser.id}, \texttt{enduser.name}) & UC4 \\
\hline
RFN-11 & Esposizione di \emph{exemplars} per le metriche Histogram (puntatore al \texttt{TraceId}) & UC4 \\
\hline
\end{tabularx}
\end{table}

\begin{table}[h]
\caption{Tabella del tracciamento dei requisiti qualitativi}
\label{tab:requisiti-qualitativi}
\begin{tabularx}{\textwidth}{lXl}
\hline\hline
\textbf{Requisito} & \textbf{Descrizione} & \textbf{Use Case}\\
\hline
RQN-1 & La telemetria non deve degradare sensibilmente le prestazioni in produzione & - \\
\hline
RQN-2 & Le metriche non devono generare un numero eccessivo di serie temporali in Prometheus (cardinalità controllata) & - \\
\hline
RQD-1 & Le dashboard devono privilegiare la visualizzazione della latenza (percentili) rispetto al solo \emph{throughput} & UC2 \\
\hline
\end{tabularx}
\end{table}

\begin{table}[h]
\caption{Tabella del tracciamento dei requisiti di vincolo}
\label{tab:requisiti-vincolo}
\begin{tabularx}{\textwidth}{lXl}
\hline\hline
\textbf{Requisito} & \textbf{Descrizione} & \textbf{Use Case}\\
\hline
RVN-1 & I segnali prodotti devono essere compatibili con qualsiasi backend \gls{otel}-compatibile & - \\
\hline
RVN-2 & Le impostazioni devono essere modificabili per ambiente (sviluppo, \emph{staging}, produzione) tramite \emph{appsettings} & - \\
\hline
\end{tabularx}
\end{table}

\chapter{Diagrammi delle classi del sistema di telemetria}
\label{app:diagrammi-classi}

In questa appendice sono raccolti i diagrammi delle classi del sistema di
telemetria descritti, in forma narrativa, nel capitolo~\ref{cap:progettazione}.
I diagrammi adottano uno stile \gls{uml} semplificato: rettangoli con nome e
membri salienti, frecce tratteggiate per le dipendenze d'uso, freccia con
triangolo vuoto per la relazione di implementazione di interfaccia.

\section{Pipeline condiviso in \texttt{BaseApplicationBuilder}}
\label{app:diag-pipeline}

La figura~\ref{fig:diag-pipeline} mostra il pipeline \gls{otel} centralizzato:
i tre provider \texttt{TracerProvider}, \texttt{MeterProvider} e
\texttt{LoggerProvider} sono costruiti una sola volta da
\texttt{BaseApplicationBuilder}, che vi registra rispettivamente il
\texttt{SiaParentBasedSampler}, il \texttt{PerSourceFilterProcessor} e il
\texttt{BaggageActivityProcessor} per le \emph{trace}, il filtro
\emph{exemplar} \texttt{TraceBased} per le metriche, e il \emph{sink} OTLP di
Serilog per i \emph{log}. Tutti i segnali condividono lo stesso
\texttt{ResourceBuilder}, garantendo la coerenza degli attributi
\texttt{service.name}, \texttt{sia.product\_type}, \texttt{sia.customer\_name}
e l'esportazione verso il \emph{Collector} tramite \texttt{OtlpExporter}.

\begin{figure}[h]
\centering
\begin{tikzpicture}[
    node distance=6mm and 10mm,
    every node/.style={font=\footnotesize},
    builder/.style={
        rectangle split, rectangle split parts=2,
        draw, thick, rounded corners=1pt,
        rectangle split draw splits=true,
        align=left, inner sep=3pt, text width=58mm,
        fill=yellow!10
    },
    provider/.style={
        rectangle split, rectangle split parts=2,
        draw, thick, rounded corners=1pt,
        rectangle split draw splits=true,
        align=left, inner sep=3pt, text width=42mm,
        fill=blue!5
    },
    comp/.style={
        rectangle, draw, thick, rounded corners=1pt,
        align=center, inner sep=3pt, text width=34mm,
        fill=orange!10
    },
    dep/.style={-{Stealth[length=2.5mm]}, thick, dashed},
    use/.style={-{Stealth[length=2.5mm]}, thick}
]

\node[builder] (builder) {
    \textbf{BaseApplicationBuilder}
    \nodepart{second}
    \texttt{+ ConfigureBuilder()}\\
    \quad registra i 3 provider OTel\\
    \quad costruisce ResourceBuilder\\
    \quad registra middleware enduser
};

\node[provider, below left=12mm and 18mm of builder] (tp) {
    \textbf{TracerProvider}
    \nodepart{second}
    \texttt{+ Sampler}\\
    \texttt{+ Processors[]}\\
    \texttt{+ OtlpExporter}
};
\node[provider, below=12mm of builder] (mp) {
    \textbf{MeterProvider}
    \nodepart{second}
    \texttt{+ ExemplarFilter = TraceBased}\\
    \texttt{+ OtlpExporter}
};
\node[provider, below right=12mm and 18mm of builder] (lp) {
    \textbf{LoggerProvider (Serilog)}
    \nodepart{second}
    \texttt{+ OTLP sink}\\
    \texttt{+ TraceId / SpanId enricher}
};

\node[comp, below=6mm of tp] (sampler)   {\textbf{SiaParentBasedSampler}};
\node[comp, below=6mm of sampler] (proc) {\textbf{PerSourceFilterProcessor}\\ \textbf{BaggageActivityProcessor}};

\node[comp, left=10mm of tp]  (resource) {\textbf{ResourceBuilder}\\ service.name, product\_type,\\ customer\_name};

\draw[dep] (builder.south) -- (tp.north);
\draw[dep] (builder.south) -- (mp.north);
\draw[dep] (builder.south) -- (lp.north);
\draw[use] (tp.south) -- (sampler.north);
\draw[use] (sampler.south) -- (proc.north);
\draw[dep] (resource.east) -- node[above, font=\scriptsize]{condiviso} (tp.west);
\draw[dep] (resource.east) -- (mp.west);
\draw[dep] (resource.east) -- (lp.west);

\end{tikzpicture}
\caption{Pipeline OpenTelemetry centralizzato in \texttt{BaseApplicationBuilder}: i tre provider (trace, metriche, log) condividono lo stesso \texttt{ResourceBuilder} e utilizzano \emph{sampler} e \emph{processor} custom per le \emph{trace}.}
\label{fig:diag-pipeline}
\end{figure}

\section{Pattern \texttt{Sia\{Modulo\}Diagnostics} e Adapter}
\label{app:diag-diagnostics}

La figura~\ref{fig:diag-diagnostics-pattern} illustra in forma generica la
classe \texttt{Sia\{Modulo\}Diagnostics} che ogni modulo del framework espone.
La classe è l'unico punto di istanziazione degli strumenti \gls{otel} del
modulo (pattern \emph{Static Factory}); inoltre, per i moduli che espongono
un controller, ospita un \texttt{ControllerDiagnosticsAdapter} annidato che
implementa il contratto condiviso \texttt{IControllerDiagnostics} (pattern
\emph{Adapter}). Il controller chiama il metodo condiviso
\texttt{BaseApiController.RecordMetricError} passando il singleton
\texttt{AsControllerDiagnostics}, senza conoscere l'implementazione
concreta.

\begin{figure}[h]
\centering
\begin{tikzpicture}[
    node distance=6mm and 14mm,
    every node/.style={font=\footnotesize},
    iface/.style={
        rectangle split, rectangle split parts=2,
        draw, thick, rounded corners=1pt,
        rectangle split draw splits=true,
        align=left, inner sep=3pt, text width=50mm,
        fill=green!5
    },
    cls/.style={
        rectangle split, rectangle split parts=2,
        draw, thick, rounded corners=1pt,
        rectangle split draw splits=true,
        align=left, inner sep=3pt, text width=60mm,
        fill=blue!3
    },
    diag/.style={
        rectangle split, rectangle split parts=2,
        draw, thick, rounded corners=1pt,
        rectangle split draw splits=true,
        align=left, inner sep=3pt, text width=66mm,
        fill=orange!10
    },
    nested/.style={
        rectangle split, rectangle split parts=2,
        draw, thick, rounded corners=1pt,
        rectangle split draw splits=true,
        align=left, inner sep=3pt, text width=56mm,
        fill=orange!4
    },
    impl/.style={-{Triangle[length=3mm, width=3mm, open]}, thick},
    dep/.style={-{Stealth[length=2.5mm]}, thick, dashed}
]

\node[iface] (i) {
    \textbf{\textlangle\textlangle{}interface\textrangle\textrangle{} IControllerDiagnostics}
    \nodepart{second}
    \texttt{+ ControllerErrors : Counter<long>}
};

\node[cls, right=18mm of i] (base) {
    \textbf{BaseApiController}
    \nodepart{second}
    \texttt{\# RecordMetricError(}\\
    \quad \texttt{IControllerDiagnostics,}\\
    \quad \texttt{Exception, operation)}
};

\node[diag, below=14mm of i] (diag) {
    \textbf{Sia\{Modulo\}Diagnostics}
    \nodepart{second}
    \texttt{+ Source : ActivitySource}\\
    \texttt{+ Meter : Meter}\\
    \texttt{+ OperationCount : Counter<long>}\\
    \texttt{+ OperationDuration : Histogram<double>}\\
    \texttt{+ ControllerErrors : Counter<long>}\\
    \texttt{+ AsControllerDiagnostics : IControllerDiagnostics}
};

\node[nested, below=6mm of diag] (adapter) {
    \textbf{ControllerDiagnosticsAdapter} \emph{(private sealed, nested)}
    \nodepart{second}
    \texttt{ControllerErrors =}\\ \quad \texttt{Sia\{Modulo\}Diagnostics.ControllerErrors}
};

\node[cls, right=18mm of diag] (ctrl) {
    \textbf{Sia\{Modulo\}Controller}
    \nodepart{second}
    \texttt{catch (Exception ex) \{}\\
    \quad \texttt{RecordMetricError(}\\
    \quad\quad \texttt{Diagnostics.AsControllerDiagnostics,}\\
    \quad\quad \texttt{ex, nameof(Op));}\\
    \texttt{\}}
};

\draw[impl] (adapter.north) -- node[right, font=\scriptsize]{\emph{implements}} (diag.south);
\draw[impl] (adapter.west) .. controls +(-2,0) and +(-0.5,-2) .. node[left, font=\scriptsize]{\emph{implements}} (i.south);
\draw[dep] (ctrl.north) -- node[right, font=\scriptsize]{extends} (base.south);
\draw[dep] (ctrl.west)  -- node[above, font=\scriptsize]{\emph{usa AsControllerDiagnostics}} (diag.east);
\draw[dep] (base.south) -- node[right, font=\scriptsize]{dipende da}  (i.east);

\end{tikzpicture}
\caption{Struttura generica del pattern \texttt{Sia\{Modulo\}Diagnostics} con \texttt{ControllerDiagnosticsAdapter} annidato. Il controller invoca \texttt{RecordMetricError} del metodo condiviso \texttt{BaseApiController}, passando il singleton \texttt{AsControllerDiagnostics}; l'Adapter riconcilia il \texttt{Counter} per-modulo con il contratto comune \texttt{IControllerDiagnostics}.}
\label{fig:diag-diagnostics-pattern}
\end{figure}

\section{Arricchimento del contesto utente}
\label{app:diag-enduser}

La figura~\ref{fig:diag-enduser} mostra le collaborazioni che realizzano
l'arricchimento automatico di ogni \texttt{Activity} con i \emph{tag}
\texttt{enduser.id} e \texttt{enduser.name}. Il
\texttt{EnrichActivityWithUserMiddleware} legge le \emph{claim} da
\texttt{HttpContext.User} e pubblica le voci nel \texttt{Baggage}; il
\texttt{BaggageActivityProcessor}, registrato come \texttt{ActivityProcessor}
nel \texttt{TracerProvider}, copia ogni \emph{entry} di \emph{baggage} con
prefisso \texttt{enduser.} come \emph{tag} sull'\texttt{Activity} nel proprio
\emph{hook} \texttt{OnStart}. L'effetto è che qualsiasi \texttt{Activity}
creata durante la richiesta --- sia dalle classi
\texttt{Sia\{Modulo\}Diagnostics}, sia dalle instrumentazioni automatiche di
ASP.NET Core e \texttt{SqlClient} --- riceve i \emph{tag} utente senza
modifiche al codice del modulo.

\begin{figure}[h]
\centering
\begin{tikzpicture}[
    node distance=5mm and 12mm,
    every node/.style={font=\footnotesize},
    cls/.style={
        rectangle split, rectangle split parts=2,
        draw, thick, rounded corners=1pt,
        rectangle split draw splits=true,
        align=left, inner sep=3pt, text width=52mm,
        fill=blue!3
    },
    comp/.style={
        rectangle, draw, thick, rounded corners=1pt,
        align=center, inner sep=3pt, text width=34mm,
        fill=orange!10
    },
    use/.style={-{Stealth[length=2.5mm]}, thick},
    dep/.style={-{Stealth[length=2.5mm]}, thick, dashed}
]

\node[cls] (mw) {
    \textbf{EnrichActivityWithUserMiddleware}
    \nodepart{second}
    \texttt{+ InvokeAsync(HttpContext)}\\
    \quad legge claim \texttt{userId}, \texttt{userName}\\
    \quad \texttt{Baggage.SetBaggage("enduser.id", \ldots)}
};

\node[comp, right=of mw] (bag) {\textbf{Baggage}\\ (\emph{OpenTelemetry context})};

\node[cls, below=of bag] (proc) {
    \textbf{BaggageActivityProcessor} : \textbf{ActivityProcessor}
    \nodepart{second}
    \texttt{+ OnStart(Activity)}\\
    \quad legge \texttt{Baggage.Current}\\
    \quad copia voci \texttt{enduser.*} su \texttt{Activity.SetTag}
};

\node[comp, right=of proc] (act) {\textbf{Activity}\\ (\emph{span in corso})};

\node[cls, below=of mw] (diag) {
    \textbf{Sia\{Modulo\}Diagnostics}
    \nodepart{second}
    \texttt{Source.StartActivity("sia.xxx.op")}
};

\draw[use] (mw.east) -- node[above, font=\scriptsize]{scrive} (bag.west);
\draw[use] (bag.south) -- node[right, font=\scriptsize]{legge} (proc.north);
\draw[use] (proc.east) -- node[above, font=\scriptsize]{\texttt{SetTag}} (act.west);
\draw[dep] (diag.east) -- node[above, font=\scriptsize]{crea} (proc.west);

\end{tikzpicture}
\caption{Flusso di arricchimento automatico delle \texttt{Activity} con l'identità utente. Il \emph{middleware} pubblica \texttt{enduser.id} e \texttt{enduser.name} nel \texttt{Baggage} di OpenTelemetry; il \texttt{BaggageActivityProcessor} li copia come \emph{tag} su ogni \texttt{Activity} creata, comprese quelle prodotte dalle classi \texttt{Sia\{Modulo\}Diagnostics}.}
\label{fig:diag-enduser}
\end{figure}

\chapter{Listati di codice rappresentativi}
\label{app:listati}

In questa appendice sono raccolti i frammenti di codice più estesi citati in
forma di rinvio nei capitoli del corpo principale. I frammenti sono pensati
come riferimento canonico: ogni nuovo modulo del \emph{framework} li adotta
con minime variazioni di \emph{naming}.

\section{Pattern Adapter applicato alla diagnostica del modulo CRUD}
\label{app:listing-adapter}

Il listato~\ref{lst:adapter-diagnostics-app} mostra la classe statica
\texttt{Sia\{Modulo\}Diagnostics} per il modulo CRUD, con il
\texttt{ControllerDiagnosticsAdapter} annidato che riconcilia il
\texttt{Counter} statico del modulo con il contratto comune
\texttt{IControllerDiagnostics}. Il listato~\ref{lst:recordmetricerror-call-app}
mostra il punto di chiamata, ridotto a una riga nel \emph{controller}
derivato.

\begin{lstlisting}[language={[Sharp]C},
  caption={Pattern Adapter applicato alla diagnostica del modulo CRUD},
  label=lst:adapter-diagnostics-app]
public static class SiaDevextremeCrudDiagnostics
{
    public static readonly Meter Meter =
        new("Sia.Devextreme.Crud", "1.0.0");

    public static readonly Counter<long> ControllerErrors =
        Meter.CreateCounter<long>(
            "sia.devextreme.crud.controller.errors");

    public static readonly IControllerDiagnostics
        AsControllerDiagnostics = new ControllerDiagnosticsAdapter();

    private sealed class ControllerDiagnosticsAdapter
        : IControllerDiagnostics
    {
        public Counter<long> ControllerErrors =>
            SiaDevextremeCrudDiagnostics.ControllerErrors;
    }
}
\end{lstlisting}

\begin{lstlisting}[language={[Sharp]C},
  caption={Invocazione di \texttt{RecordMetricError} dal controller derivato},
  label=lst:recordmetricerror-call-app]
catch (Exception ex)
{
    RecordMetricError(
        SiaDevextremeCrudDiagnostics.AsControllerDiagnostics,
        ex,
        nameof(GetData));
    throw;
}
\end{lstlisting}
